% 02-abbreviations.tex - Abbreviations, Definitions, and Conventions

\chapter{Document Conventions}
\label{sec:abbreviations}

\section{Abbreviations and Acronyms}

\begin{longtable}{l p{10cm}}
\toprule
\textbf{Abbreviation} & \textbf{Definition} \\
\midrule
\endfirsthead

\toprule
\textbf{Abbreviation} & \textbf{Definition} \\
\midrule
\endhead

\bottomrule
\endfoot

AVX     & Advanced Vector Extensions \\
COR     & Correlator (data format) \\
DR      & Data Recorder \\
DROS    & Data Recorder Operating Software \\
DRX     & Digital Receiver (4-bit data format) \\
DRX8    & Digital Receiver 8-bit (data format) \\
DRSU    & Data Recorder Storage Unit \\
EXT4    & Fourth Extended File System \\
FFT     & Fast Fourier Transform \\
FFTW    & Fastest Fourier Transform in the West (FFT library) \\
GbE     & Gigabit Ethernet \\
ICD     & Interface Control Document \\
LWA     & Long Wavelength Array \\
MCS     & Monitor and Control System \\
MIB     & Management Information Base \\
MJD     & Modified Julian Day \\
MPM     & Milliseconds Past Midnight \\
NDP     & Next-generation Digital Processor \\
NTP     & Network Time Protocol \\
RAID    & Redundant Array of Independent Disks \\
SPC     & Spectrometer Command \\
SSE     & Streaming SIMD Extensions \\
TBS     & Transient Buffer --- Streaming (data format) \\
TBT     & Transient Buffer --- Triggered (data format) \\
UDP     & User Datagram Protocol \\

\end{longtable}

\section{Command Parameter Types}

The following parameter types are used to describe command arguments and MIB
response fields throughout this document:

\begin{description}
    \item[\textbf{uint8}:] Unsigned integer, 8~bits wide.

    \item[\textbf{uint16}:] Unsigned integer, 16~bits wide.

    \item[\textbf{uint32}:] Unsigned integer, 32~bits wide.

    \item[\textbf{uint64}:] Unsigned integer, 64~bits wide.

    \item[\textbf{ASCII-\textit{XXX}-\#}:] An ASCII string exactly \textit{XXX}
        characters in length which is interpreted as a number.  Valid characters
        are digits and right-padding spaces only.

    \item[\textbf{ASCII-\textit{XXX}-A}:] An ASCII string exactly \textit{XXX}
        characters in length which is interpreted as a text string.  Unless
        otherwise noted, valid characters are letters, numbers, the underscore
        character, and periods.
\end{description}

\section{Mark-up Conventions}

\begin{longtable}{l p{4.5cm} p{6cm}}
\toprule
\textbf{Symbol} & \textbf{Meaning} & \textbf{Example} \\
\midrule
\endfirsthead

\toprule
\textbf{Symbol} & \textbf{Meaning} & \textbf{Example} \\
\midrule
\endhead

\bottomrule
\endfoot

\textit{italics}
    & Variable or parameter name
    & \textit{Start MPM} \\[4pt]

\texttt{fixed-width}
    & Literal parameter or format value.
      A single-quote (\texttt{'}) in a
      literal format denotes a space.
    & \texttt{AB'\_} = ``A'' + ``B'' + space + ``\_'' \\[4pt]

\texttt{<\ldots>}
    & Variable substitution; brackets
      omitted from actual format
    & \texttt{A<B>\_} = ``A'' + value of \textit{B} + ``\_'' \\[4pt]

\texttt{[\ldots]}
    & Optional parameter; brackets
      omitted from actual format
    & \texttt{A[B]\_} = ``A'' + optional \textit{B} + ``\_'' \\

\end{longtable}

\section{Numeric Representation Convention}

Numbers, units, and their associated prefixes and suffixes conform to the
standard of IEC~60027-2.  Specifically:

\begin{itemize}
    \item Prefixes Ki, Mi, Gi, and Ti refer to $2^{10}$, $2^{20}$, $2^{30}$,
          and $2^{40}$, respectively.
    \item Prefixes K, M, G, and T refer to $10^{3}$, $10^{6}$, $10^{9}$,
          and $10^{12}$, respectively.
    \item For binary size or rate: uppercase B represents a byte, lowercase b
          indicates a bit (e.g., MB = megabyte = $1{,}000{,}000$~bytes;
          Kib = kibibit = $1{,}024$~bits).
\end{itemize}
